\section{The scanner estimate from its inputs}
\label{sec:al_scanner}

Section~9 of \cite{OpenAI2026AreaLaw} scans the collar of a finite target in
rounds to select a charge round and interpolation parameter with small defect
energy (Proposition~9.2). Fix exponents $0<\ell<1$, $0<\kappa<\mu<1-\ell$,
$0<e<1$ with $\kappa\le e$, and $\nu>0$, let $W\ge1$ be the total metric weight
(possibly depending on $n$), and set $L=\lfloor n^{1-\ell}\rfloor$,
$D=\lceil n^\kappa\rceil$, $m=\lfloor n^\mu\rfloor$, $\epsilon=n^{-\nu}$, and
$\delta_n=C\epsilon^{-1}(n^{e-\mu}+a^{1/4}(\log n)^{C})$. The statement below
takes as hypotheses at scale $n$ the conclusions of Lemma~9.1 (scan histories),
Proposition~7.4 (log-norm evolution), Proposition~8.1 (comparators), Lemma~2.1
(entropy continuity), and Lemma~3.1 (marginal tails); it is Proposition~9.2
restricted to those conclusions rather than the unconditional source
proposition. The finite leaf measures are required to have mass $1/2$ only for
$0<p<1$ (with no mass equality at an inactive terminal leaf), the logarithmic
norm is continuous on $[0,1]$ with matching values at $0$ and $1$, and the
charge defect is almost everywhere measurable on $(0,1)$ (hence integrable on
$[0,1]$ by its nonnegative upper bound).

% Source: OpenAI2026AreaLaw, sections/08-scanner.tex, lines 32--41 (the scales);
% the three comparison bounds are the estimates used at lines 469, 480, 519--520.
\begin{lemma}[Numerical scale estimates]
    \label{lem:al_scanner_scale_facts}
    \lean{TNLean.PEPS.AreaLaw.Scan.eventually_scaleFacts}
    \leanok
    Fix $0<\ell<1$, $0<\kappa<\mu<1-\ell$, $\nu>0$, and $0<e<1$ with
    $\kappa\le e$, and set $L=\lfloor n^{1-\ell}\rfloor$,
    $m=\lfloor n^\mu\rfloor$, $K=\lfloor L/(8m)\rfloor$,
    $D=\lceil n^\kappa\rceil$, and $\epsilon=n^{-\nu}$
    \cite[Section~9]{OpenAI2026AreaLaw}. There is $C_1\ge1$, depending only on
    the exponents, such that for all sufficiently large $n\in\mathbb N$,
    $n\ge2$, $K,m,D\ge1$, $L\le n$, $0<\epsilon<1$, and
    \begin{align}
        \frac{n^e}{m}   & \le C_1n^{e-\mu},  &
        D               & \le C_1n^e,        &
        \frac{nD^2}{mK} & \le C_1n^\ell D^2.
        \label{eq:al_scanner_scale_facts}
    \end{align}
\end{lemma}
\begin{proof}\leanok
    For large $n$, $n^\mu,n^{1-\ell}\ge2$ and $n^{1-\ell-\mu}\ge48$. Since
    $\lfloor x\rfloor\ge x/2$ for $x\ge2$, $m\ge n^\mu/2$ and
    $L\ge n^{1-\ell}/2\ge24m$, so $K\ge1$ and $mK\ge L/16\ge n^{1-\ell}/32$;
    combining with $\kappa\le e$ and $\lceil x\rceil\le2x$ ($x\ge1$) gives~(\ref{eq:al_scanner_scale_facts})
    with $C_1=32$.
\end{proof}

\begin{lemma}[Continuity cost at trace distance $n^{-500}$]
    \label{lem:al_continuity_cost}
    \lean{TNLean.PEPS.AreaLaw.Scan.eventually_continuityCost_le}
    \leanok
    For any $C,\lambda\ge0$, all sufficiently large $n\in\mathbb N$ and all
    $t\le n^{-500}$ satisfy \cite[Lemma~2.1; Proposition~9.2]{OpenAI2026AreaLaw}
    \begin{align}
        C\sqrt{t}\bigl(1+Cn^2\lambda\bigr)\le n^{-247}.
        \label{eq:al_continuity_cost}
    \end{align}
\end{lemma}
\begin{proof}\leanok
    For $n\ge\max\{1,C(1+C\lambda)\}$, $\sqrt t\le n^{-250}$ and
    $1+Cn^2\lambda\le(1+C\lambda)n^2$, so
    $C\sqrt t(1+Cn^2\lambda)\le n^{-250}C(1+C\lambda)n^2\le n^{-247}$.
\end{proof}

\begin{lemma}[Continuity cost with a fixed distance coefficient]
    \label{lem:al_continuity_fixed_coefficient}
    \lean{TNLean.PEPS.AreaLaw.Scan.eventually_continuityCost_le_of_fixed_coefficient}
    \leanok
    For any $A,C,\lambda\ge0$ (with $\sqrt t=0$ for $t<0$), all sufficiently
    large $n\in\mathbb N$ and all $t\le An^{-500}$ satisfy
    \cite[Proposition~9.2]{OpenAI2026AreaLaw}
    \begin{align}
        C\sqrt{t}\bigl(1+Cn^2\lambda\bigr)\le n^{-247}.
        \label{eq:al_continuity_fixed_coefficient}
    \end{align}
\end{lemma}
\begin{proof}\leanok
    \uses{lem:al_continuity_cost}
    Set $B=1+A>0$, $\overline C=C\sqrt B\ge C$, and $u=t/B\le n^{-500}$. Then
    $\sqrt t=\sqrt B\sqrt u$ for all $t\in\mathbb R$, so
    Lemma~\ref{lem:al_continuity_cost} applied to $(\overline C,\lambda)$ gives
    \begin{align}
        C\sqrt t\bigl(1+Cn^2\lambda\bigr)
        =\overline C\sqrt u\bigl(1+Cn^2\lambda\bigr)
        \le\overline C\sqrt u\bigl(1+\overline Cn^2\lambda\bigr)
        \le n^{-247}.
        \notag
    \end{align}
\end{proof}

\begin{lemma}[Growth of the rounded band count]
    \label{thm:area_law_band_count_growth}
    \lean{TNLean.PEPS.AreaLaw.Scan.ScannerExponents.eventually_le_K}
    \leanok
    For the scanner exponents above and
    $K_n=\lfloor\lfloor n^{1-\ell}\rfloor/(8\lfloor n^\mu\rfloor)\rfloor$
    \cite[Section~9]{OpenAI2026AreaLaw}, every integer $B\ge0$ satisfies
    $B\le K_n$ for all sufficiently large $n$.
\end{lemma}
\begin{proof}\leanok
    Since $1-\ell-\mu>0$, eventually $n^{1-\ell-\mu}\ge16B$ and
    $n^{1-\ell},n^\mu\ge2$, so
    \begin{align}
        8B\lfloor n^\mu\rfloor
        \le8Bn^\mu
        \le\tfrac12 n^{1-\ell}
        \le\lfloor n^{1-\ell}\rfloor,
        \notag
    \end{align}
    and dividing by $8\lfloor n^\mu\rfloor>0$ yields $B\le K_n$.
\end{proof}

\begin{lemma}[Scanner estimate from the conclusions of its inputs]
    \label{thm:al_scanner_from_inputs}
    \lean{TNLean.PEPS.AreaLaw.Scan.ScanData,
        TNLean.PEPS.AreaLaw.Scan.ScanData.SXt_le,
        TNLean.PEPS.AreaLaw.Scan.ScanData.SQt_le,
        TNLean.PEPS.AreaLaw.Scan.ScanData.terminalBound,
        TNLean.PEPS.AreaLaw.Scan.ScanData.terminalBound_le_of,
        TNLean.PEPS.AreaLaw.Scan.ScanData.IntegratedChargeBound,
        TNLean.PEPS.AreaLaw.Scan.ScanData.chargeDefect_le_choiceGainSum,
        TNLean.PEPS.AreaLaw.Scan.ScanData.chargeDefect_le,
        TNLean.PEPS.AreaLaw.Scan.ScanData.goodOld_le,
        TNLean.PEPS.AreaLaw.Scan.exists_budgets,
        TNLean.PEPS.AreaLaw.Scan.ScanData.intervalIntegrable_chargeDefect,
        TNLean.PEPS.AreaLaw.Scan.ScanData.integratedChargeBound,
        TNLean.PEPS.AreaLaw.Scan.exists_selected_density,
        TNLean.PEPS.AreaLaw.Scan.ScanData.energySum_le,
        TNLean.PEPS.AreaLaw.Scan.exists_defectEnergy_le,
        TNLean.PEPS.AreaLaw.Scan.scanner_estimate}
    \leanok
    There are constants $C$, independent of $W$ and depending only on the
    exponents and input constants, such that for all sufficiently large $n$ and
    every scan at scale $n$ satisfying the conclusions of Lemmas~9.1, 2.1, 3.1
    and Propositions~7.4, 8.1 (with leaf-measure mass identities, charge-defect
    measurability, log-norm continuity/differentiability and derivative bounds,
    initial norm bound, finite leaf measures, and matching adjacent endpoint
    norms for $k\ge0$, and normalized final comparison for $k\ge1$), we have
    $B_{\mathrm{sh}}\le C(nL^e+n)$, $B_{\mathrm{exc}}\le CnD$,
    $\mathcal B\le CnD(\log n)^{C}$, and a uniform remainder $\rho_k\to0$ such
    that for every $k\ge1$ some charge round and $p_k\in[\epsilon/2,\epsilon]$
    satisfy $\mathcal Q(p_k)/(KnD)\le\delta_n+\rho_k$ and
    \begin{align}
        E_{\mathrm{def},k}\le Cn^\ell D^2W^2(\log n)^{C}
        \bigl(\delta_n^{1/8}+\epsilon+n^{-100}\bigr)+Cn^{-1000}+\rho_k.
        \label{eq:al_scanner_energy}
    \end{align}
    Moreover, if $C_e$ is the entropy-input constant and
    $\widetilde S_X,\widetilde S_Q(j)$ are the target and shell entropies of the
    truncated state, then for every $j\le L$,
    \begin{align}
        \widetilde S_X    & \le C_e n^{1+e}+n^{-247}, &
        \widetilde S_Q(j) & \le C_e(nL^e+n)+n^{-247}.
        \label{eq:al_scanner_truncated_entropies}
    \end{align}
\end{lemma}
\begin{proof}\leanok
    \uses{lem:al_scanner_scale_facts,lem:al_continuity_fixed_coefficient}
    Applying Lemma~\ref{lem:al_continuity_fixed_coefficient} with
    $A=C=C_{\rm in}\ge1$ and $\lambda=\log q\ge0$ to $t\le C_{\rm in}n^{-500}$
    bounds the continuity differences by $n^{-247}$, giving~(\ref{eq:al_scanner_truncated_entropies})
    and, with $\mathcal B\le CnD(\log n)^{12}$, the bounds on
    $B_{\mathrm{sh}},B_{\mathrm{exc}},\mathcal B$.
    For $k\ge1$, integrating the interior log-norm derivative inequality over
    $[0,1]$, retaining the nonnegative charge on $[\epsilon/2,\epsilon]$, and
    telescoping across rounds gives, for the charge rounds $T$ ($|T|=nm$) and
    $c_k=KnD(\delta_n+\rho_k)$,
    \begin{align}
        \sum_{r\in T}\int_{\epsilon/2}^{\epsilon}\mathcal Q_{r,k}(p)\,dp
        \le |T|\frac{\epsilon}{2}c_k.
        \label{eq:al_scanner_charge_sum}
    \end{align}
    Since $T\ne\varnothing$,~(\ref{eq:al_scanner_charge_sum}) yields $r\in T$
    with $\int_{\epsilon/2}^{\epsilon}\mathcal Q_{r,k}(p)\,dp\le(\epsilon/2)c_k$;
    if $\mathcal Q_{r,k}>c_k$ on $[\epsilon/2,\epsilon]$, then
    \begin{align}
        0<\int_{\epsilon/2}^{\epsilon}\bigl(\mathcal Q_{r,k}(p)-c_k\bigr)\,dp
        =\int_{\epsilon/2}^{\epsilon}\mathcal Q_{r,k}(p)\,dp-\frac{\epsilon}{2}c_k\le0,
        \notag
    \end{align}
    so some $p_k\in[\epsilon/2,\epsilon]$ satisfies
    $\mathcal Q_{r,k}(p_k)/(KnD)\le\delta_n+\rho_k$. At $p_k$, H\"older's
    inequality with exponents $(8/7,8)$ bounds old good leaf energies by
    $N(\mathcal Q/N)^{1/8}$, new and bad leaves contribute at most
    $(\epsilon+n^{-200})$ times total weight, and the gap of $\overline H_k$
    gives~(\ref{eq:al_scanner_energy}).
\end{proof}

\begin{theorem}[Scanner estimate for the collar scan from the evolution and comparator inputs]
    \label{thm:al_scanner_actual_scan}
    \lean{TNLean.PEPS.AreaLaw.Scan.ScanInputs,
        TNLean.PEPS.AreaLaw.Scan.HistoryFacts.mono,
        TNLean.PEPS.AreaLaw.Scan.CollarScan.scanner_estimate_actualScan}
    \leanok
    \uses{thm:al_scanner_from_inputs,thm:al_scan_actual_history_facts,
        thm:al_scan_truncation_inputs}
    Fix scanner exponents, $R$, $J\ge0$, $\Delta>0$, and $C_T\ge0$. There are
    $C_r$ and, for each $q$, constants $C_0,c_0,C_1>0$ such that for all scanner
    constants with dimension $q$, normalized gap $g$, common constant $\ge C_0$,
    logarithmic power $\ge4$, and sampling constant $\le c_0$, the following
    holds for all sufficiently large $n$: for any gapped ground vector $\Omega$
    on $\Lambda$ and collar scan at the scales of
    \cite[Section~9]{OpenAI2026AreaLaw} with $r_0=\lceil C_r(\log n)^2\rceil$
    and charge-list constant $C_1$ under the test geometry and entropy inputs of
    \cite[Proposition~9.2]{OpenAI2026AreaLaw}, there is a ground vector
    $\widetilde\Omega$ of $\widetilde H$ with gap $g/2$ satisfying
    Lemma~\ref{thm:al_scan_truncation_inputs} such that, for every auxiliary
    dimension, total weight $W\ge1$ ($a=W/K$), and replica vectors, if the scan
    rounds satisfy the conclusions of Propositions~7.4 and~8.1 for
    $(W,\widetilde\Omega)$ with the replica conditions of
    Lemma~\ref{thm:al_scanner_from_inputs}, then the conclusions of
    Lemma~\ref{thm:al_scanner_from_inputs} hold (Proposition~9.2 conditional
    on Propositions~7.4 and~8.1 for the scan rounds).
\end{theorem}
\begin{proof}\leanok
    \uses{thm:al_scanner_from_inputs,thm:al_scan_actual_history_facts,
        thm:al_scan_truncation_inputs}
    Combine the truncation inputs for $\widetilde\Omega$ from
    Lemma~\ref{thm:al_scan_truncation_inputs} and the scan-round history
    inputs from Lemma~\ref{thm:al_scan_actual_history_facts} (which persist
    under increasing $C_0$ and the logarithmic power once $\log n\ge1$ and
    decreasing $c_0$) with the assumed conclusions of Propositions~7.4 and~8.1,
    and apply Lemma~\ref{thm:al_scanner_from_inputs}.
\end{proof}
